
\subsubsection{Neural Theorem Proving}

LeanDojo (Yang et al., 2023) represents current state-of-the-art in neural theorem proving, achieving 48.9\% success rate on held-out theorems from the Lean mathematical library using a retrieval-augmented ByT5-based transformer. The system employs best-first search with LLM-based tactic suggestion. GPT-f (Polu \& Sutskever, 2020) and PACT (Han et al., 2022) demonstrated that transformer models can learn effective proof strategies from formal proof corpora.

These systems focus on improving tactic suggestion quality rather than resource allocation decisions. They inherit fixed timeout strategies from symbolic theorem provers.

\subsubsection{Timeout Strategies in Automated Reasoning}

Traditional automated theorem provers employ non-adaptive timeout strategies. Z3 (de Moura \& Bjørner, 2008) and Vampire (Kovács \& Voronkov, 2013) use resource limits based on structural metrics like clause counts and proof depth. These heuristics are theorem-agnostic and do not leverage semantic information available in neural provers.

\subsubsection{Uncertainty Quantification and OOD Detection}

Confidence calibration methods (Guo et al., 2017) aim to align predicted probabilities with empirical frequencies. Out-of-distribution detection techniques (Hendrycks \& Gimpel, 2017; Liang et al., 2018) use model uncertainty to identify inputs from unseen distributions. However, these methods have not been systematically applied to sequential decision problems in neural reasoning systems.

\subsubsection{Symbolic Divergence Detection}

Symbolic theorem provers use cycle detection via state hashing (Nieuwenhuis \& Rubio, 1995) and proof state growth metrics to identify potential non-termination. These structural signals are orthogonal to neural confidence measures.

\subsubsection{Positioning}

No prior work has combined LLM confidence analysis with symbolic signals for termination detection in neural theorem proving. The present work proposes this combination and claims to validate it empirically, though as noted above, the validation uses synthetic rather than real data.
